\chapter{结论与展望}

\section{结论}
本文围绕“基于 3DGS 与 StyleGAN2 的 MRI 医学图像拟生成平台的设计与实现”开展工作。课题最终完成了 MRI 数据预处理、二维切片生成、连续切片组织、点云初始化、3DGS 训练结果接入和前后端平台展示。整个工作不是单独复现某个算法，而是把多个原本分散的步骤整理成一条可以运行、可以展示、可以复核的工程链路。

具体来看，平台已经实现四项主要功能。第一，系统可以读取 NIfTI 格式 MRI 数据，并完成切片提取、灰度归一化、尺寸调整、低信息切片过滤和训练数据整理。第二，项目基于 StyleGAN2-ADA 完成二维 MRI 切片生成模型训练，并把训练权重接入平台推理流程。用户可以通过随机种子和截断参数生成可复现的切片结果。第三，项目把连续切片序列转换为 PLY 初始点云，并构建规则相机参数，使 3DGS 可以在该输入上完成三维表达实验。第四，平台采用 FastAPI 和 React 实现前后端分离，整合结果概览、切片生成、医学影像查看、任务管理、三维展示和 AI 助手等页面。

本课题的创新点主要体现在工程适配和流程集成上。StyleGAN2、StyleGAN2-ADA 和 3DGS 的基础算法来自已有公开研究，本文不把算法原理本身作为原创贡献。本文的工作在于把 StyleGAN2 的二维生成能力接入 MRI 切片流程，把连续切片组织为点云和 3DGS 输入，并把训练、推理、结果管理和可视化统一到 FastAPI + React 平台中。这样处理后，算法结果不再只停留在命令行和离线文件中，而是可以通过页面查看、解释和复核。

实验结果表明，平台已经具备教学演示和实验复核所需的基本能力。StyleGAN2 训练链路形成了可加载权重，平台可以基于该权重生成 MRI 切片。3DGS 部分在 180 张连续切片监督下完成训练，并输出三维高斯模型、渲染结果和训练指标。第 30000 轮训练达到 L1 为 0.0459、PSNR 为 20.0909 dB，完整点云包含 556424 个高斯点，前端预览点云包含 60000 个点。这些结果说明，平台已经把二维生成、三维表达和结果展示连接成可操作流程。

本文也明确限定了平台的使用边界。当前系统主要用于教学演示、毕业设计展示和工程实验验证，不用于临床诊断。实验数据以脑部 T1 MRI 为主，数据类型和病例来源仍然有限。3DGS 部分使用规则切片相机和单序列灰度图像，能够形成可展示的三维表达，但还不能完整代表真实医学空间结构。评价方面，本文使用 L1、PSNR 和可视化对比进行分析，尚未加入医学专家评价、下游任务验证和更全面的医学可信性评估。

\section{展望}
后续工作可以先从三维表达继续完善。当前系统使用规则相机参数完成 3DGS 训练，能够支撑工程展示，但还没有充分利用 NIfTI 中的真实体素间距、层厚和空间方向信息。后续可以把这些信息加入相机参数构建过程，并尝试层间连续性约束、结构平滑约束和高斯点裁剪策略，让三维表达更加接近 MRI 体数据本身的空间关系。

评价体系也需要进一步补充。当前论文主要使用 L1、PSNR 和可视化对比来说明训练过程。后续可以加入 SSIM、FID、LPIPS、分割一致性或下游任务性能等指标。对于医学图像生成任务，还需要医学专业人员参与主观评价，从结构合理性、解剖一致性和教学可用性等角度检查生成结果。

平台工程方面也还有提升空间。当前系统已经能够通过本地 FastAPI 服务和 React 前端完成主要功能展示，但运行时仍依赖本地模型路径、数据路径、GPU 环境和外部接口配置。后续可以通过配置文件、环境变量、任务队列和容器化部署规范运行环境。对于耗时较长的训练任务，还可以加入更完善的日志追踪、失败恢复和资源监控功能。

应用场景方面，本文主要围绕脑部 T1 MRI 切片展开。后续可以扩展到多序列 MRI、CT 或其他公开医学图像数据集，也可以研究跨模态生成、数据增强和三维教学展示等任务。若未来能够获得更充分的数据、专家评价和伦理审查支持，平台可以继续作为医学图像生成方法比较、三维重建流程教学和医学人工智能实验课程的辅助工具。
